\documentclass[10pt,a4paper]{amsart}
\usepackage[margin=19mm]{geometry}
\usepackage{amsmath,amssymb,mathtools}
\usepackage[scr=rsfs,cal=euler]{mathalfa}
\usepackage{xfrac}
\usepackage[unicode,hidelinks,bookmarks=false]{hyperref}
\newcommand{\ie}{i.e.\ }
\newcommand{\cf}{cf.\ }
\newcommand{\resp}{respectively\ }
\newcommand{\Subsection}{\subsection}
\providecommand{\nobreakdash}{\nobreak\hbox{-}}
\setlength{\emergencystretch}{2em}
\multlinegap0pt
\usepackage{bm}
\usepackage{bbm}
\usepackage{dsfont}
\usepackage{xspace}
\usepackage{cancel}
\usepackage{mathtools}
\usepackage{amsthm}
\makeatletter
\@ifundefined{theorem}{\newtheorem{theorem}{Theorem}}{}
\@ifundefined{lemma}{\newtheorem{lemma}[theorem]{Lemma}}{}
\makeatother
\title[On Locally Generalized radical linear groups over rings]{On Locally Generalized radical linear groups over rings}
\author{Le Van Chua, Bui Xuan Hai}
\date{Source: arXiv:2608.15136v1 (CC BY 4.0)}
\begin{document}
\maketitle

\noindent\emph{Source and licence.} On Locally Generalized radical linear groups over rings. Version arXiv:2608.15136v1 (CC BY 4.0), distributed under the Creative Commons Attribution 4.0 International licence, as recorded in the arXiv licence metadata for this exact version and shown on the abstract page https://arxiv.org/abs/2608.15136v1. Full rights notice: licenses/arxiv-2608.15136-CC-BY-4.0.txt.\par\smallskip

\noindent\textit{Editorial scope.} Counted unit: the Lemma whose source label is \texttt{lem:5.1} in the article (source0.tex lines 386--392, source label \texttt{lem:5.1}), delivered together with the article's own complete original proof of it (source0.tex lines 393--415, one \texttt{proof} environment of 23 lines). No mathematics is written, repaired or re-derived: statement and proof are the article's own text line for line. Required context retained uncounted: none. Every definition, display and label of those lines is retained unchanged. Omitted from this package and not redistributed: 1--385, 416--904. Nothing inside a retained excerpt is omitted or altered: every retained body is delivered exactly as the article's lines are written, so no omission receipt of any kind accompanies this package. Citations: the retained excerpts carry no citation command at all, so no bibliography entry of the article is transcribed and no reference is redistributed. Reference adaptation: none was needed and none was made; every reference inside the delivered unit resolves inside it, so no unresolved reference or citation is produced. Printed slips kept literal: no printed word, symbol, hypothesis or number of the counted statement or of its proof was altered. Layout and class: the article's own class is \texttt{amsart} with packages including amsmath, amssymb, graphicx, tikz, color, enumerate, hyperref; the wrapper is the lane's pinned \texttt{amsart} editorial wrapper at 10pt on A4, which restates the theorem-like environments the retained text needs and adds the article's own macro definitions that the retained text needs (none), with extra wrapper packages (bm, bbm, dsfont, xspace, cancel, mathtools). The delivered numbering is the unit's own. Assets: no retained span contains a figure environment, a graphics-inclusion command, a PDF-page inclusion, a table or a TikZ picture; no image file is copied into this package, the package declares no asset dependency and no image file is redistributed with it. Licence: version arXiv:2608.15136v1 of this article is distributed under the Creative Commons Attribution 4.0 International licence, as recorded in the arXiv licence metadata for this exact version and shown on the abstract page https://arxiv.org/abs/2608.15136v1. A full rights notice is delivered at licenses/arxiv-2608.15136-CC-BY-4.0.txt. Characters: every retained excerpt is pure ASCII, so the delivered engine, XeLaTeX with its default Latin Modern text font, reports no missing character.
\begin{lemma}\label{lem:5.1} Let $R$ be a ring with its Jacobson radical $J=J(R)$, $R^\times$ the unit group of $R$, and $I$ an ideal of $R$. Then, the following assertions hold:
\begin{itemize}
\item [(1)] $G(I)$ is a submonoid of $R^\times.$
\item [(2)] If $I\subseteq J,$ then $G(I)$ is a normal subgroup of $R^\times$ and $R^\times/G(I)\cong (R/I)^\times.$ In particular, $G(J)$ is normal in $R^\times$ and $R^\times/G(J)\cong (R/J)^\times.$ Moreover, the unit group $R^\times$ has a normal series $$1\trianglelefteq G(I)\trianglelefteq G(J)\trianglelefteq R^\times.$$
\item [(3)] If $I$ is a nilpotent ideal of nilpotency class $k$, then $G(I)$ is a nilpotent group of nilpotency class at most $k.$ Moreover, $G(I)\trianglelefteq R^\times$ and $$R^\times/G(I)\cong (R/I)^\times.$$
\end{itemize}	
\end{lemma} 

\begin{proof} (1) For every $x=1+a, y=1+b\in G(I)$, where $a,b\in I$, we have
 $$xy=(1+a)(1+b)=1+(a+b+ab)\in G(I),$$	
which shows that $G(I)$ is closed under multiplication, and the identity element $1\in G(I).$ Consequently, $G(I)$ is a submonoid of the unit group $R^\times.$

(2) Suppose that $I\subseteq J$. If $a\in I,$ then $a\in J,$ and so $1+a\in R^\times.$ Thus, $G(I)\subseteq R^\times.$ Hence, by conclusion of (1), we get that $G(I)$ is a subgroup of $R^\times$. Further, for any $r\in R^\times$ and $a\in I$, we have
$$rxr^{-1}=r(1+a)r^{-1}=1+rar^{-1}\in G(I),$$ which shows that $G(I)$ is normal in $R^\times.$ Since $I\subseteq J,$ it follows that $G(I)\subseteq G(J).$ As a consequence, we get the normal series 
	$$1\trianglelefteq G(I)\trianglelefteq G(J)\trianglelefteq R^\times.$$
Now, consider any non-zero element $a\in R$. Observe that $a\in R^\times$ if and only if $\overline{a}=a+I\in (R/I)^\times$ because $I\subseteq J$. It follows that the natural homomorphism $R^\times\longrightarrow (R/I)^\times$ given by $a\mapsto \overline{a},$ is a surjective homomorphism whose kernel is $G(I),$ and hence $R^\times/G(I)\cong (R/I)^\times.$
		
(3) Suppose that $I$ is a nilpotent ideal of nilpotency $k$. For all $r\in \{1,2,\ldots,k\},$ let $$G_r(I)=1+I^r\subseteq G(I).$$
Since $I^r\subseteq J$, in view of (2), $G_r(I)$ is a normal subgroup of $R^\times$, and we have the following  normal series of subgroups
$$1=G_k(I)\le G_{k-1}(I)\le\cdots\le G_1(I)=G(I).$$
We shall prove that this is a central series of $G(I)$. Indeed, for each $1\leq s, r\leq k,$ let $x=1+a\in G_s(I),$ with $a\in I^s$ and $y=1+b\in G_r(I),$ with $b\in I^r.$ We have
	\begin{gather}
		\begin{split}
			[x,y]-1&=x^{-1}y^{-1}xy-1\\
			&=x^{-1}y^{-1}(xy-yx)\\
			&=x^{-1}y^{-1}(ab-ba)\\
			&\in R\cdot(I^s\cdot I^r+I^r\cdot I^s)\subseteq I^{s+r}.
		\end{split}\notag
	\end{gather}
	This implies that $[x,y]\in G_{s+r}(I),$ and so $[G_s(I),G_r(I)]\le G_{s+r}(I).$ Hence, we conclude that $G(I)$ is a nilpotent group of nilpotency class at most $k$, and at this point,  the proof of the lemma is  complete.
\end{proof}

\end{document}
